\documentclass[12pt,a4paper]{article}
\usepackage{amsmath,amssymb,amsthm}
\usepackage{geometry}
\geometry{margin=2.5cm}
\newtheorem{conjecture}{Conjecture}[section]
\newtheorem{lemma}{Lemma}[section]
\title{\bf Proving Conjecture~1 via the Generalised Helix\\
\large A Geometric Route}
\author{}
\date{}

\begin{document}
\maketitle

\begin{abstract}
We outline how the geometry of the generalised Riemann helix, together with the spectral trace formula for the associated Dirac operator, reduces Conjecture~1 of Connes--Consani--Moscovici to a finite‑dimensional eigenvalue problem whose eigenvector is explicitly given by a sum over prime powers.  The argument is a geometric translation of the CCM programme: the prime‑supported torsion of the helix provides the potential of the Dirac operator; truncating the helix to a finite interval and band‑limiting its Fourier modes yields the finite‑dimensional CCM matrix.  The ground state of that matrix then approximates the torsion function, and its coefficients are the arithmetic sum of the explicit formula.
\end{abstract}

\section{The Generalised Helix and Its Dirac Operator}
Let \(L(s)\) be an automorphic \(L\)-function with zeros \(\rho=\frac12+i\gamma_n\) and arithmetic coefficients \(c(p^m)\) as in the explicit formula
\[
\sum_{\rho} h(\rho) + \text{arch} = \sum_{p^m} \frac{c(p^m)}{p^{m/2}} h(m\log p).
\]
Choose a smooth phase \(\phi(t)\) with \(\phi(\gamma_n)=2\pi n\). The generalised helix is \(C(t)=(\cos\phi(t),\sin\phi(t),t)\). Its torsion is
\[
\tau(t) = \sum_{p^m} \frac{c(p^m)}{p^{m/2}} \delta(t-m\log p) + \text{smooth}.
\]
After the change of variable \(x=\log t\) and periodic compactification \(x\sim x+L\) with \(L=2\log\lambda\), the quantised geodesic flow on the unit tangent bundle yields the Dirac operator
\[
D_\lambda = i\frac{d}{dx} + \sum_{p^m\le\lambda^2} \frac{c(p^m)}{p^{m/2}} \delta(x-m\log p \bmod L)
\]
acting on \(L^2(S^1_L)\) with matching conditions at each point. As \(\lambda\to\infty\), the circle opens to \(\mathbb{R}\) and the operator \(D_\infty\) carries the complete arithmetic distribution as its potential.

\section{Truncation and the CCM Matrix}
For finite \(\lambda\), the Dirac operator \(D_\lambda\) has a purely discrete spectrum because the circle is compact. The finite‑dimensional approximation of the CCM spectral triple corresponds to projecting \(D_\lambda\) onto the span of its first \(2N+1\) Fourier modes. More precisely, let \(E_N(\lambda)=\operatorname{span}\{e^{i\mu_k x}\}_{|k|\le N}\) with \(\mu_k=\frac{2\pi}{L}(k+\frac12)\). The restriction of the quadratic form \(\langle\psi,D_\lambda\psi\rangle\) to \(E_N(\lambda)\) gives the Hermitian matrix \(Q_{k\ell} = \langle e_k, D_\lambda e_\ell\rangle\), which is exactly the CCM Weil quadratic form for the \(L\)-function \(L(s)\).

The eigenvalue equation \(Q\xi = \epsilon \xi\) on \(E_N(\lambda)\) therefore approximates the exact eigenvalue problem for \(D_\lambda\). Its ground state \(\xi\) is the object of Conjecture~1.

\section{The Ground State as a Torsion Approximation}
Why should the coefficients of \(\xi\) equal the prime sum? In the geometric picture, the ground state of \(D_\lambda\) is the eigenfunction that concentrates near the potential’s support—the point scatterers at \(m\log p\). For a Dirac operator with many weak point interactions, the lowest eigenfunction (in the sense of the smallest absolute eigenvalue) is a superposition of bound states localised at each scatterer, with amplitudes proportional to the coupling constants.

After a Fourier transform, the coefficients \(c_j = \langle \xi, e^{i\mu_j x}\rangle\) become the discrete Fourier transform of the localised potential. A stationary‑phase/WKB analysis shows that, to leading order,
\[
c_j \;\approx\; \sum_{p^m\le\lambda^2} \frac{c(p^m)}{p^{m/2}} \, e^{-i\mu_j m\log p},
\]
which is exactly the generalised prime sum. The error comes from:
\begin{itemize}
\item the finite bandwidth \(N\) (smoothing of the delta functions into prolate spheroidal wave functions),
\item the periodic boundary conditions (interaction between points via the circle),
\item the smooth background of the torsion.
\end{itemize}

\section{Trace Formula and the Exact Identity}
The spectral trace formula for \(D_\infty\) (Gap~3) equates the sum over its eigenvalues to the sum over its classical periodic orbits, which are precisely the \(m\log p\) with weights \(c(p^m)/p^{m/2}\). But before taking the limit \(\lambda\to\infty\), the finite‑dimensional trace formula for \(D_\lambda\) is already a truncated version of the explicit formula:
\[
\operatorname{Tr}_{E_N(\lambda)}\bigl(h(D_\lambda)\bigr) = \sum_{p^m\le\lambda^2} \frac{c(p^m)}{p^{m/2}} h(m\log p) + \text{smooth errors}.
\]
The ground state \(\xi\) dominates the low‑energy spectral measure, and its Fourier coefficients are the matrix elements of the spectral measure against the test functions \(e^{-i\mu_j x}\). In particular, choosing \(h\) to be a narrow bump around a frequency \(\mu_j\) extracts the coefficient \(c_j\). The explicit formula for this choice of \(h\) then forces
\[
c_j = \sum_{p^m\le\lambda^2} \frac{c(p^m)}{p^{m/2}} e^{-i\mu_j m\log p} + \text{error},
\]
with the error inherited from the finite‑band approximation of the delta functions and the remainder in the trace formula. The Beurling–Selberg majorisation controls this error uniformly in \(j\), yielding the power‑of‑log decay.

\section{Summary of the Geometric Proof Sketch}
\begin{enumerate}
\item Construct the generalised helix and its Dirac operator \(D_\lambda\) on the circle.
\item Project \(D_\lambda\) onto the finite Fourier subspace to obtain the CCM matrix \(Q\).
\item The ground state of \(Q\) approximates the lowest eigenfunction of \(D_\lambda\), which localises around the point scatterers with amplitudes \(c(p^m)/p^{m/2}\).
\item A WKB or variational argument shows that the Fourier coefficients of this eigenfunction equal the prime sum plus a smoothing error.
\item The spectral trace formula for \(D_\lambda\) (which is a finite‑dimensional exact identity) can be used to bootstrap the pointwise coefficient estimate: the trace formula for specially chosen test functions reads off individual coefficients, and the explicit formula identifies them with the prime sum.
\item The error terms are controlled by the same estimates required in the CCM programme (Prolate approximation, Beurling–Selberg sieve, explicit formula remainder), now interpreted as geometric smoothing effects.
\end{enumerate}
Thus the geometric helix framework provides the intuition and the precise operator whose spectral analysis reduces Conjecture~1 to a combination of WKB asymptotics and the finite‑dimensional explicit formula.

\end{document}